% Auto-generated by scripts/06_emit_structure.py
% Source: scans 182–187
% Section id: appendix/3
% Phase 3 deliverable. Footnotes (Phase 4), citations (Phase 5),
% and Wade-Giles → Pinyin (Phase 6) are not yet applied here.

\chapter{The Size of the Sample}
\label{app:3}

% source: scan 182, printed 165
The Size of the Sample

\footnote[1]{In what follows, as throughout the book, I attempt to distinguish between fragment (a fragment of bone or shell) and piece (a bone or shell published as one rubbing, but which may consist of one or more fragments). Yi-pien 7758, for example, is a fragment; Ping-pien 45 (formed by joining Yi-pien 4864 and 7758) is a piece. That the various counts cited in n. 2 do not always observe this distinction reduces their accuracy, but not, I believe, significantly. On this point, cf. Tung (1953), p. 375. On the way duplicates affect total fragment counts, see, for example, Tung (1953), p. 378; (1954a), p. 120; cf. n. 2 below. The problem of duplicates is a most serious one, particularly for some collections. Of the 655 rubbings in Chien-shou (1917), for example, only 60 are not reproduced in Hsü-pien (1933); of the 1,125 rubbings in Fu-yin (1925), very few are not reproduced in Hsü-pien (Mickel, letters of 26 May 1973 and 5 August 1976). For indexes to duplicates, see ch. 5, n. 52, above.}% phase4: unanchored note 1


\footnote[2]{Hu Hou-hsüan (1951), p. 65, and (1952), p. 6, estimated that 161,189 fragments (or pieces?) had been found. This was dismissed as a wild exaggeration by Tung (1953), pp. 380, 381, who revised Hu's figures downward to arrive at a total of 54,113 unpublished pieces and 42,005 published pieces, making 96,118 fragments (or pieces? in all; this figure he subsequently raised to 109.617 (Tung [1954a]. p. 120; [1965], p. 13). Ch'en Meng-chia (1956, p. 47, arrived at an estimated total of 98,000 (cf. his figures, ibid., pp. 655-657). Complete accuracy in such matters would not justify the efforts involved, and most scholars now speak of approximately 100,000 fragments found (e.g., Chang Ping-ch'üan [1967], p. 831 Kaizuka [1967], p. 205). To that figure of 100,000 must be added the 21 scapulas excavated in December 1971 (Kuo Mo-jo [1972]) and the 7,150 bone and shell fragments excavated in 1973 (KK [1975.1], p. 37). This means that about 107,000 fragments have now been found and recorded. Note, however, Mickel's estimate that he has now made index cards for over 10,000 cases of duplicate publication (letter of 24 August 1976); this suggests that the figure of 107,000 may eventually prove too large. 165 U U D}% phase4: unanchored note 2

% source: scan 183, printed 166
As we have seen (appendix 2, sec. 2), it took, on the average, 16 bone fragments to make one
scapula, and 35 shell fragments to make one plastron. Before we attempt to translate our 107,000
fragments into numbers of whole scapulas and plastrons, however, it is necessary to estimate the
number of blank fragments which would have formed part of the original corpus but which,
lacking inscriptions, tended to be discarded as valueless.\footnote[3]{Blanks may be of two types: (1) fragments of scapulas or plastrons which, though used for divination, were never inscribed with writing; or (2) blank fragments from scapulas or plastrons whose divination inscriptions did not extend to this area of the bone or shell. Inability to distinguish between the two types renders our calculations imprecise. The suggestion of Liu Yüan-lin (1974), p. 128, relying on physical criteria (cf. ch. 4, n. 150), that the pre-Wu Ting oracle bones at An-yang were not inscribed is plausible but not yet proven.}3 In the case of the approximately 32,000
fragments excavated scientifically between 1928 and 1973, no adjustments are necessary since the
count presumably reflects the total number of fragments found, blank as well as inscribed.\footnote[4]{The Academia Sinica excavated 24,918 fragments between 1928 and 1937 (Tung [1953], p. 378). Twenty-one scapulas were excavated in 1971 and 7,150 bone and shell fragments in 1973 (see n. 2). This makes a total of 32,089. On the way blanks may or may not have been registered in the first dig, see ch. 4, n. 184.}4 A
study of these fragments indicates that approximately 48 percent of all fragments excavated were
blank. If we apply that figure to the approximately 75,000 fragments collected privately, we may
conclude that they must have been accompanied originally by approximately 69,000 fragments
that were presumably blank and not preserved.
Adding these 144,000 75,000 + 69,000) private fragments to the 32,000 excavated scientifically, we may estimate that 176,000 fragments, both blank and inscribed, must have been
excavated in the An-yang region between 1899 and\footnote[6]{This figure, of course, ignores the thousands of fragments that must have been thrown away, ground up for medicine or fertilizer, or lost in countless other ways over the centuries (cf. Hu Hou-hsüan [1944a], p. 4b; [1955], p. 11). For example, nearly fifty thousand fragments stored in Menzies' house were reported destroyed (or buried?) in 1928 (Yetts [1933], p. 673); these may have been part of the store of over ten thousand fragments dug up on the campus of Chi-nan University in 1952, most of which had rotted away (Ming-hou, preface, p. 3). There is no way to estimate the total numbers lost, but since they are lost they have no bearing on this part of our enquiry, save to remind us that a certain proportion of the Shang corpus that was once found is now, thanks to the activities of man, forever irretrievable.} 1973.6
If we assume (following the argument in appendix 2) that plastrons and scapulas were used
in approximately equal quantities and that it took, on the average, 16 bone fragments to make
one scapula and 35 shell fragments to make one plastron, then our 176,000 fragments may be taken
as equivalent to 3,450 scapulas and 3,450 plastrons. These figures are, of course, gross approximations, but they are suggestive."

\footnote[5]{Of the 24,918 fragments excavated by the Academia Sinica in twelve digs between 1928 and 1937, approximately 13,000 pieces (52 percent) were published in Chia-pien and Yi-pien (Tung [1953], p. 378; cf. Ch’en Meng-chia [1956], p. 655). An identical ratio may be obtained by counting the number of published and unpublished (and presumably blank) fragments in all pits containing more than fifty fragments, namely, A26, E9, F36, and F37 from the first dig; TLK and longitudinal trench 13 ping (north) from the third dig; E16 and H20 from the fifth dig; B119, YH6, YH127, YH330, and YH344 from the thirteenth dig; and the 1973 finds. The statistics may be taken from Tung (1929b); (1929c); (1949); (1949a); Shih Chang-ju (1959), pp. 320-323; KK (1975.1), p. 37. Previous estimates of the proportion of blank to inscribed fragments have ranged from over 90 percent (Li Chi [1933], p. 575) to approximately 50 percent (Hu Hou-hsüan [1944a], pp. 4b-5a) to over 50 percent (Tung [1953], p. 375).}% phase4: unanchored note 5


\footnote[7]{Hu Hou-hsüan (1944a), pp. 4a-5b, arrives at far larger totals: 16,003 plastrons and 11,858 scapulas; and if the 1971 and 1973 finds (n. 2) were included, his figures would be still larger. His estimate of fragments found, however, appears inflated (n. 2), and his view that it took only ten fragments to make a plastron and five to make a scapula is unsupported by the evidence (see appendix 2, sec. 2).}% phase4: unanchored note 7

% source: scan 184, printed 167
<!-- cjk-tess-overlay: 了 一 -->

If, for example, we arbitrarily assume that the Shang kings used I scapula and 1 plastron a
day, the 3,450 scapulas and 3,450 plastrons would have lasted less than 10 years. For the present
materials to have lasted the approximately 150 years from Wu Ting's accession to the end of the
dynasty, the Shang kings could have used only I scapula and I plastron every 16 days.”
Such a rate of use is impossibly low. We know of divinations about the night made on a
daily basis; we know the fortune of the coming week was divined every ten days; we know of
cases where five plastrons were used in approximately eleven days (sec. 3.7.4.1). Presumably
these were not the only divinations made on those days; these were not the only oracle bones used.
The fact that carving was done on a mass-production basis (sec. 2.9.3); the fact that several hundred
shells might be received in one shipment, and that 30 or 40 or even 100 shells might be ritually
prepared at one time (sec. 1.4)-all this suggests that the diviners used more than I bone or I
shell every 16 days. If we temporarily, and conservatively, assume that the Shang cracked and
inscribed 1 scapula and 1 plastron a day, then, at the most, our 176,000 fragments represent only
1/16 (6.25 percent) of the original corpus."
C C C
2.
Estimate on the Basis of Marginal Notations
The proportion of the original inscriptions that we now possess may also be estimated
from a study of the marginal notations. For example, the notation "Ch'üeh sent in 250 (shells)"
has now been found on some 22 plastrons or plastron fragments from pit YH127.11 If we assume
that every plastron sent in at that time was inscribed with the same notation, 12 we are still missing
228 plastrons (or, more accurately, plastron fragments) with that notation. Similarly, 6 plastrons
or plastron fragments from the same pit contain the notation "Wo¹³ brought in (?) 1,000 (shells)";
presumably, the other 994 notations have not yet been discovered. 14 By contrast, Ping-pien 13.9
8. Presumably, the rate of use varied from
period to period, topic to topic, diviner to diviner.
In period I, a single plastron might be used up
in one day (Ping-pien 8; fig. 8), or a set of five
plastrons might be used up in eleven days (see
ch. 3, n. 119); in period III, a set of five plastrons
might be used up in twenty-six days (Chia-pien
3917); in period V. a set of five scapulas might
be used for at least seventy days (Hsü-ts'un 2.972
[D]; fig. 10). More research is needed to show if
there was a trend in the later periods to use
scapulas and plastrons over longer periods of time.
9. For the 150 years of my revised short
chronology, see appendix 4, sec. 6; if the long
chronology of 273 years is used, the Shang would
have divined I scapula and 1 plastron every 29
days.
10. This figure is based on the revised short
chronology (see appendix 4, sec. 6); if the long
chronology is used, the figure falls to less than

\section[percent. And if the Shang used more than one]{\origsecnum{3.5}percent. And if the Shang used more than one}
\label{sec:appendices-app03:3.5}

bone and shell a day, the percentages would be
still lower.
II. Chui-ho 127; Ping-pien 196 (Cr); 355 (Cr);
357; 359 (Cr); 369 (Cr); 374 (Cr); 376; 389 (Cr);
539; 579 (Cr); 586 (Cr); 616 (Cr); Yi-pien 754
(Cr); 3289 (Cr); 3300; 4420; 4541; 7050; 7123
(Cr); 7153; 7491 (Cr); cf. Yi-pien 5256.1-2 (Itō
[1962a], pp. 254-255). For the turtle identifications, see tables 3 and 4.
12. The assumption is unverifiable (cf. Shirakawa [1957], p. 61, n. 1); some of the shells from
this shipment may not have been so inscribed.
On the other hand, the possibility that Ch'üeh
sent in two or more sets of 250 shells during period
I cannot be excluded. The facts that the plastrons
with this notation were all excavated from YH127
and that, to the extent that we can measure and
identify them, many were of similar size (cf. Tung
[1949a], pp. 268-269) and of the same species
(Chinemys reevesi) (n. 11) make it plausible to
assume, at least in this case, that they all belonged
to the same shipment.
13. See ch. 1, n. 72.
14. The six pieces are Yi-pien 1053 (Cr); 2684
(Cr); 2702 (= Ping-pien 268); 3432 (Cr); 6686
(Cr); 6967 (Cr) (S352.2). Two more fragments,
Yi-pien 909 and 6860 (S352.2), may belong to the
same shipment, but their notations are incomplete.
These plastrons were all found in pit YH127.
% source: scan 185, printed 168
<!-- cjk-tess-overlay: 喜 -->

and 21.9, with the marginal notation “Chih Fei (?) sent in 2 (shells)," represent a case where the
2 shells sent in at one time have presumably been found, starting and ending the same set of 5
plastrons (for this set, see sec. 3.7).
It is possible to quantify such data by considering the number of shells supplied by Ch'üeh,
Chu, Chu, and Hsi, who were among the major suppliers (S256.1-256.2). The figures in
table 36 support the logical conclusion that the larger the number of pieces sent in, the smaller the
percentage we are now likely to have¹5 and that the percentage falls sharply as the number of pieces
sent in rises over 10. That is, the average percentage recovered for groups of 10 plastrons or less is
about 57 percent; for groups in the range from 11 to 99 plastrons, the average percentage recovered
is about 6 percent; for groups of 100 plastrons or more it is about 3.5 percent. Again, working
with crude approximations, it appears that the median size of all income notations was 10 units;
the average size was 88.16 This suggests—and again I would stress the roughness of these figures--
that while we may have approximately 57 percent of the more numerous, smaller shipments, we
may have only some 2 to 6 percent of the larger ones. And the bulk of the Shang receipts appears
to have come in these larger shipments. As a rough estimate, therefore, it would seem plausible
to suppose that for period I at least, when these notations were being recorded (sec. 4.3.1.6), we
have no more than 5 percent of the materials that were sent in to the Shang. 17
3.
Estimate by Scheduled Divinations
A third approach for estimating the degree to which our present store of inscriptions
represents the original corpus takes advantage of the fact that under the New School of periods
IIb and V the Shang divined about certain topics according to a rigid and continuous schedule
(for examples see appendix 5, sec. 2). If we assume that Tsu Chia reigned for 20 years, 18 we would
expect to find, according to Tung's study of the five rituals in period IIb, divinations about 23
sacrificial cycles, or, since there were 245 divinations in each cycle, 5,635 ritual divinations. 19
In fact, Shima's Sōrui (sec. 3.3.2) records about 327 inscriptions, 5.8 percent of the number expected.
15. The records of ritual preparation (ch. I,
n. 71) confirm that the smaller our original
sample, the more fragments we are likely to have
(this approach was first used by Tung [1954],
pp. 460, 463). Thus, the socket notation "Fu
Ching ritually prepared so many pairs (of scapulas)" (S151.3) is found on 9 sockets; these 9
notations refer to a total of 62 bones (duplicate
cases being excluded to the best of my ability);
the similar notation for Fu Hao is found on 7
sockets, referring to a total of 41 bones (\$151.3-4);
for Yie, it is found on 19 sockets, referring to
104 bones (S152.1). In these cases, therefore, we
have 14.5 percent, 17 percent, and 18 percent of
the original number of bones (always assuming,
of course, that all the sockets were so inscribed in
each case), but the average number of bones
involved in these 3 cases is less than 6.
16. These figures are derived from Hu Houhsüan (1944a), pp. 11b-12b. His tables are out
of date and incomplete, but there is no reason to
think that they seriously distort the true situation.
17. It is not certain that the Shang used and
carved inscriptions on all the materials that were
sent in. It is conceivable that when Ch'üeh sent
in 250 shells, the Shang discarded, or used for
other purposes, all but twenty-two of them. But
this seems unlikely. I assume, pending evidence
to the contrary, that bones and shells were valued
and their entry noted primarily because they were
needed for Shang divination, and that they were
so used.
18. For the twenty-year reign length, see
appendix 4, sec. 5.
19. Five rituals offered to each of 29 ancestors
and 20 ancestresses call for a total of 245 inscriptions. For the source of these figures, see n. 20.
% source: scan 186, printed 169
<!-- cjk-tess-overlay: 一 -->

If, following the long chronology, Tsu Chia is accorded a reign of 33 years, involving 38 sacrifice
cycles, our present inscription sample represents only 3.5 percent of the original. 20
Similarly, if we assume that Ti Yi and Ti Hsin reigned for 50 years, 21 and divined about 50
sacrificial cycles, involving, at this point, 285 divinations in each cycle, 22 we now have only 587
five-ritual inscriptions (4 percent) out of the total of 14,250 that we believe were originally divined.
If we use the long chronology, that figure would be reduced to 2 percent. 23 If we count the "ten-day
plus appended sacrifice" inscriptions from period V, we find a similar situation. Assuming a 50-year
period, we now have only 257 (14 percent) out of an expected 1,800 divinations. If we use the long
chronology, that figure would be reduced to 7 percent. 24
These figures need to be revised slightly upward by taking account of inscriptions, such as
those in Menzies and White, not included in Shima's Sörui. 25 But again, this method suggests that
if the short chronology applies, we possess little more than from 4 to 14 percent of the inscriptions
originally made; and if the long chronology applies, from 2 to 7 percent.
4.
Conclusions: Rate of Use
It is true that numbers of pieces (or fragments) (discussed in sec. 1 and 2 above) and
numbers of inscriptions (discussed in sec. 3) are not fully commensurable since one piece or fragment
may contain several inscriptions. It is also true that the representative nature of the sample may
vary by topic and period. Nevertheless, all approaches support the conclusion that, on the average,
the number of inscribed fragments we now possess is, at most, no more than about 5 to 10 percent
of those originally in existence.
If the 3,450 scapulas and 3,450 plastrons to which we reduced our present corpus (sec. 1)
represent, let us say, 5 percent of the materials which the Shang actually used, the total inscription
corpus would originally have been recorded on some 69,000 scapulas and 69,000 plastrons. If we
use the short chronology, which allots approximately 150 years to the Wu Ting-Ti Hsin era, the
Shang diviners would have used approximately 1.3 scapulas and 1.3 plastrons a day, or the bones
20. This approach was pioneered by Tung
(1945), pt. 1, ch. 3, pp. 16a-b; pt. 2, ch. 8, p. 1a,
who found only 61 divinations (0.6 percent) out
of an expected 9,310: here I revise his figures by
counting the inscriptions listed by S91.3-4;
127.2-4; 250.2-3: 398.3-4; 492.4-493.2. These
calculations assume, of course, that the New
School reforms were instituted in the first year of
Tsu Chia's reign; if instituted later, the percentage
figures would rise accordingly. For the long and
short chronologies, see appendix 4.
21. For the fifty-year reign length, see appendix
4, sec. 5.
22. Five rituals offered to each of 33 ancestors
and 24 ancestresses call for a total of 285 inscriptions. For the source of these figures, see n. 23
below.
23. Following the approach of Tung (1945),
pt. 1, ch. 3, pp. 18b, 20b, but revising his figures
by counting the inscriptions listed at S92.1-3;
128.1-3;250.4-251.2;398.4-399.2;493.4-494.4.
24. Following the approach of Tung (1945),
pt. 1, ch. 3, p. 2ob, but revising his figures by
counting the inscriptions listed at S166.4-168.2.
If we follow Tung by assigning Ti Hsin only 52
sacrificial cycles for his 63 years, then, using the
long chronology, we would have only 8 percent
of the original inscriptions.
25. See, for example, Hsü Chin-hsiung (1970b).
The 1973 finds are not likely to affect our calculations since they are reported to contain few
period V fragments and none from period II
(KK [1975.1], p. 39).
% source: scan 187, printed 170
and shells from some 6 bovids and 13 turtles each ten-day week.\footnote[26]{The rate of use would decrease if the long chronology of 273 years is used, or if it is argued that we now have more than 5 percent of the original corpus. For example, using the short chronology (150 years), if our present corpus is 10 percent of the original corpus, rate of use would have been about 0.63 scapulas and 0.63 plastrons a day; with the long chronology (273 years), if our present corpus is 5 percent of the original corpus, rate of use would have been 0.69 scapulas and 0.69 plastrons a day; if 10 percent, rate of use would have been 0.35 scapulas and 0.35 plastrons a day. The improbable lowness of all these figures suggests that our present corpus may be less than 5 percent of the original corpus, or that the short chronology is to be preferred to the long, or both.} 26 Such an estimate, crude though
it is, is presumably of the right magnitude.\footnote[27]{The fact that we have so few complete sets extending over five bones or shells (for examples, see ch. 2, n. 50) is further indication of the small percentage of the original corpus we now possess. By period V, for example, sets appear to have been divined regularly on five separate bones or shells (as in fig. 10), but we do not have, so far as I know, a set of fragments from the five bones in even one set. 170}27
